\scanpage{19}
% archive sheet number: 18
% L1--3, upper diagrams.
\diagram{\begin{tikzcd}[ampersand replacement=\&,column sep=large,row sep=large]
\prod_x G({\unsure{\hat{\mathcal O}_x|\hat O_x}})\arrow[d] \& 0\arrow[l]\arrow[d] \& 0\arrow[d,rightsquigarrow] \& 0\arrow[l]\arrow[d,rightsquigarrow]\\
\prod_{\mathrm{local}}G(\hat K_x) \& G(K)\arrow[l] \& G\text{-{\unsure{diviseur|diviseurs}} sur }X \& G(K)\arrow[l]\\
\& \& H^1(X,G) \& H^0(X,G)
\end{tikzcd}}
% L4--9, small note below the first square.
i.e. on {\unsure{veut}} que {\unsure{ça corresponde|cela corresponde}} à un $G$-diviseur {\unsure{véritable}} sur $X$ [{\unsure{sinon donc|sinon dont}} le complexe {\unsure{continuerait-il}} ---]
% The horizontal rule between the two portions is a page separator.
\par\medskip\hrule\medskip
% L10--20: whole lower portion crossed by three oblique strokes.
\[\deleted{{\unsure{\sum|\prod|\coprod}}}\]
\[
\left\{\begin{aligned}
&\deleted{(f_x)_{x\in X}}\quad\deleted{\text{$G$-répartition}}\\
&\deleted{(\xi_D)_{D\in\mathrm{Div}\,\text{{\unsure{premiers}}}}}\qquad
\deleted{\xi_D\in J_{\underline{\mathcal O}_D/k}^*}
\end{aligned}\right.
\]
\[
\begin{aligned}
\deleted{x\in D}\qquad
\deleted{f_x\in G(K_x)}&\deleted{\longrightarrow\text{{\unsure{classe de}} }G(K_D)/\deleted{{\unsure{G(\mathcal O_D)}}}}\quad
\deleted{{\unsure{\simeq|\longrightarrow}}\ G(K_D)/G(\mathcal O_D)^{(1)}}\\
&\deleted{\hspace{4em}\simeq\operatorname{Hom}(J_{\underline{\mathcal O}_D/k}^*,G(k))}
\end{aligned}
\]
\deleted{Donc $(f_x,\xi_D)_x\in G(k)$ est défini.}

\deleted{Je dis que $\sum_{x\in|D|}(f_x,\xi_D)$ existe. En effet la question étant locale, on est {\unsure{ramené au cas}} où $(f_x)$ est formé d'unités,}
